\label{exp:ceb02}

\subsubsection*{Постановка}
\textbf{Физический вопрос.}
Позволяют ли дополнительность двугранных углов тетраэдра~$K_4$ и октаэдра $\arccos(+1/3)+\arccos(-1/3)=\pi$, диофантова соизмеримость ковалентных хребтов и двухбитный регистр основания $\mathbb{Z}_4$ вывести без свободных параметров угловую щель спиральной фрустрации $\Delta\theta_{\mathrm{frust}} \approx 7{,}36^\circ$, масштабный закон $5:6:7 = d_{\mathrm{out}}(E) : |\square| : d_{\mathrm{out}}(V)$ для трёх планетарных типов жизни ($m \in \{-1,0,+1\}$) и количественные законы генетики и мутаций (\cref{thm:bio-frustration-and-three-life-types,thm:bio-genetics-chargaff-mendel-mutations})?

\textbf{Что проверяем.}
\begin{enumerate}[label=\textbf{(\alph*)},leftmargin=*,itemsep=0.25em]
\item Выполнение тождества $\arccos(+1/d)+\arccos(-1/d) \equiv \pi$ ($70{,}53^\circ + 109{,}47^\circ = 180^\circ$) и положительность щели пяти тетраэдров вокруг ребра $\Delta\theta_{\mathrm{frust}} = 2\pi - 5\arccos(1/3) \approx 7{,}3561^\circ$.
\item Выполнение масштабного закона $5:6:7$ для длин связей $R_{\mathrm{bond}}^{(m)} \in \{\frac{10}{3},\ 4,\ \frac{14}{3}\}\caKappaGeom a_0 \approx \{1{,}247,\ 1{,}497,\ 1{,}746\}\ \text{\normalfont\AA}$ ($-\text{C{=}N}-$, $-\text{C--N}-$, $-\text{Si--N}-$ / $-\text{P--O}-$) и толщин 18-звенных трансмембранных стержней $L_{\mathrm{TM}}^{(m)} \in \{60,\ 72,\ 84\}\caKappaGeom a_0 \approx \{22{,}45,\ 26{,}94,\ 31{,}43\}\ \text{\normalfont\AA}$ (азотосома Титана, бислой Земли и мембрана Венеры).
\item Точное совпадение долей синонимических транзиций 3-й позиции кодона ($14/16 = 7/8 = 87{,}5\%$ стеблей и $60/64 = 15/16 = 93{,}75\%$ кодонов), индекса эпигенетического стебля $s(\text{CG}) = 7 = d_{\mathrm{out}}(V) \to s(\text{TG}) = 3 = d$ и отношений расщепления Менделя $(1:2:1)$ и $(3:1) = (d:1)$.
\item Воспроизведение универсального показателя метаболической аллометрии Клейбера $\alpha_{\mathrm{Kleiber}} = d/d_{\mathrm{in}}(\square) = 3/4$ ($\beta + \gamma = -1/4 + 1/4 = 0$), уставок клеточного гомеостаза трёх миров $T_{\mathrm{homeo}}^{(-1,0,+1)} = \frac{10}{81}T_{\mathrm{soft}}^{(m)} \approx \{88{,}4,\ 309{,}4,\ 618{,}8\}\ \text{К}$, квантов энерговалюты $|\Delta G_{\mathrm{vivo}}^{(-1,0,+1)}| = \frac{8}{3}E_{\mathrm{soft}}^{(m)} \approx \{15{,}9,\ 55{,}6,\ 111{,}1\}\ \text{кДж}\cdot\text{моль}^{-1}$ (с инвариантом $\frac{8/3}{10/81} = 21{,}6$), мембранных спайков $\Delta V_{\mathrm{spike}}^{(-1,0,+1)} = s_0 V_{\mathrm{soft}}^{(m)} \approx \{30{,}9,\ 108{,}0,\ 216{,}0\}\ \text{мВ}$, спектральных окон фотосинтеза $\bigl[\lambda_{13}^{(m)},\, \lambda_8^{(m)}\bigr]$ ($1{,}55\text{--}2{,}51\ \text{мкм}$ на Титане; $441{,}6\text{--}717{,}6\ \text{нм}$ на Земле; $220{,}8\text{--}358{,}8\ \text{нм}$ на Венере) и капсидных сфер Каспара--Клуга $N_{\mathrm{shell}}(n) = 10n^2+2$ (\cref{thm:bio-allometry-atp-and-kinetics,thm:bio-neuro-photosynthesis-harmony}).
\end{enumerate}

{\sloppy\textbf{Рабочая гипотеза.}
Все перечисленные структурные параметры трёх типов жизни и генетических законов воспроизводятся из констант носителя~\latGCKshort\ (\texttt{Bio\_frustration\_three\_life\_types}).\par}

\textbf{Наблюдаемые.}
$\theta_{\mathrm{tet}} + \theta_{\mathrm{oct}}$, $\Delta\theta_{\mathrm{frust}}$, $R_{\mathrm{bond}}^{(-1,0,+1)}$, $L_{\mathrm{TM}}^{(-1,0,+1)}$, $f_{\mathrm{stems}}(\mathrm{Ti}_3)$, $f_{\mathrm{codons}}(\mathrm{Ti}_3)$, $s(\text{CG})$, отношения Менделя (табл.~\ref{tab:exp-ce-b-02}).

\textbf{Критерий согласия с теорией.}
Точное выполнение угловых и комбинаторных тождеств; согласие длин связей и толщин мембран с экспериментальными и расчётными данными.

\subsubsection*{Теоретическая часть}
По теоремам~\ref{thm:bio-frustration-and-three-life-types} и~\ref{thm:bio-genetics-chargaff-mendel-mutations} геометрические размеры биополимеров в окне $m\in\{-1,0,+1\}$ определяются множителем $|\square|+m \in \{5,6,7\}$, а устойчивость кода к транзициям младшего бита $\mathrm{cod}\oplus 1$ — вырождением третьей позиции в $14$ из $16$ стеблей Румера.

\textbf{Проверка размерностей.}
Углы, доли синонимичности и стеблевые индексы безразмерны; длины связей и толщины мембран имеют размерность длины:
\[
  [\Delta\theta_{\mathrm{frust}}] = \bigl[f_{\mathrm{codons}}(\mathrm{Ti}_3)\bigr] = [s(\text{CG})] = 1,\qquad
  \bigl[R_{\mathrm{bond}}^{(m)}\bigr] = \bigl[L_{\mathrm{TM}}^{(m)}\bigr] = [\caKappaGeom a_0] = \mathrm{L}\ (\text{\normalfont\AA}).
\]

\subsubsection*{Ход эксперимента}
{\sloppy Воспроизведение: \texttt{python verify\_principles.py} \texttt{--suite sm} (\texttt{Bio\_frustration\_three\_life\_types}).\par}

\begin{enumerate}[label=\textbf{Шаг \arabic*:},leftmargin=*,itemsep=0.35em]
\item Вычислить двугранные углы $\theta_{\mathrm{tet}}, \theta_{\mathrm{oct}}$ и щель фрустрации $\Delta\theta_{\mathrm{frust}}$.
\item Вычислить $R_{\mathrm{bond}}^{(m)}$, $P_\alpha^{(m)}$, $L_{\mathrm{TM}}^{(m)}$ и $h_{\mathrm{stack}}^{(m)}$ для $m \in \{-1,0,+1\}$.
\item Проверить доли синонимических транзиций $7/8$ и $15/16$, индекс $s(\text{CG})=7$ и свести результаты в табл.~\ref{tab:exp-ce-b-02}.
\end{enumerate}

\begin{table}[htbp]
\centering
\scriptsize
\setlength{\tabcolsep}{1.5pt}
\caption{Верификация геометрической фрустрации, масштаба $5:6:7$ трёх типов жизни и законов генетики (CE-B-02)}
\label{tab:exp-ce-b-02}
\begin{tabular}{@{}llrrr@{}}
\toprule
Инвариант / наблюдаемая & Формула из носителя & Модель & Опыт / эталон & Статус \\
\midrule
Сумма углов $\theta_{\mathrm{tet}} + \theta_{\mathrm{oct}}$ & $\arccos(+1/3) + \arccos(-1/3)$ & $180{,}000^\circ$ & $\pi = 180^\circ$ & PASS \\
Щель фрустрации $\Delta\theta_{\mathrm{frust}}$ & $360^\circ - d_{\mathrm{out}}(E)\arccos(1/3)$ & $7{,}3561^\circ$ & $7{,}3561^\circ$ (Бёрдейк--Кокстер) & PASS \\
Звено $m=-1$ ($-\text{C{=}N}-$, Титан) & $\frac{10}{3}\caKappaGeom a_0$ & $1{,}247\ \text{\normalfont\AA}$ & $1{,}25\ \text{\normalfont\AA}$ (имин/нитрил) & PASS \\
Звено $m=0$ ($-\text{C--N}-$, Земля) & $\frac{12}{3}\caKappaGeom a_0 = 4\caKappaGeom a_0$ & $1{,}497\ \text{\normalfont\AA}$ & $1{,}48\ \text{\normalfont\AA}$ (пептид) & PASS \\
Звено $m=+1$ ($-\text{P--O}-$ / $-\text{Si--N}-$) & $\frac{14}{3}\caKappaGeom a_0$ ($\text{H}_3\text{PO}_4$ / Венера) & $1{,}746\ \text{\normalfont\AA}$ & $1{,}74\ \text{\normalfont\AA}$ (фосфат/силазан) & PASS \\
Стэкинг дуплекса $h_{\mathrm{stack}}^{(-1,0,+1)}$, \text{\normalfont\AA} & $\frac{3(6+m)}{2}\caKappaGeom a_0 = \{7{,}5,9{,}0,10{,}5\}\caKappaGeom a_0$ & $2{,}81;\ 3{,}37;\ 3{,}93$ & $3{,}38\ \text{\normalfont\AA}$ (B-ДНК при $m=0$) & PASS \\
Мембраны $L_{\mathrm{TM}}^{(-1,0,+1)}$, \text{\normalfont\AA} & $\{5,6,7\}\cdot 12\caKappaGeom a_0 = \{60,72,84\}\caKappaGeom a_0$ & $22{,}5;\ 26{,}9;\ 31{,}4$ & $22{,}5$ (Титан); $27{,}0$ (Земля) & PASS \\
Синонимичность $\mathrm{Ti}_3$ (стебли / кодоны) & $14/16 = 7/8;\ 60/64 = 15/16$ & $87{,}5\%;\ 93{,}75\%$ & $14/16;\ 60/64$ (код NCBI-1) & PASS \\
Отношение $\mathrm{Ti}/\mathrm{Tv}$ (фон / WGS / экзом) & $s_0 = 0{,}5;\ d-1 = 2{,}0;\ d = 3{,}0$ & $0{,}5;\ 2{,}0;\ 3{,}0$ & $0{,}5;\ 2{,}0\text{--}2{,}1;\ 3{,}0\text{--}3{,}3$ & PASS \\
Эпигенетический стебель $\text{CG}\to\text{TG}$ & $(1\ll 2)\mid 3 = d_{\mathrm{out}}(V) \to d$ & $7 \to 3$ & $\text{CpG}\to\text{TpG}$ ($5$-mC) & PASS \\
Расщепление Менделя $F_2$ (ген.\ / фен.) & $[1,1]*[1,1];\ (d : d_{\mathrm{in}}(V))$ & $(1:2:1);\ (3:1)$ & $(1:2:1);\ (3:1)$ & PASS \\
Аллометрия Клейбера $\alpha,\ \beta,\ \gamma$ & $d/d_{\mathrm{in}}(\square) = 3/4;\ \mp 1/4$ & $0{,}75;\ \mp 0{,}25$ & $0{,}75;\ \beta+\gamma=0$ & PASS \\
Гомеостаз $T_{\mathrm{homeo}}^{(-1,0,+1)}$, К & $\frac{10}{81}T_{\mathrm{soft}}^{(m)}$ ($m=-1,0,+1$) & $88{,}4;\ 309{,}4;\ 618{,}8$ & $90{,}6$ (Титан); $309{,}8$ & PASS \\
Макроэрг $|\Delta G_{\mathrm{vivo}}^{(-1,0,+1)}|$ & $\frac{8}{3}E_{\mathrm{soft}}^{(m)}$ ($\Delta G/\Theta = 21{,}6$) & $15{,}9;\ 55{,}6;\ 111{,}1$ & $54\text{--}57\ \text{кДж/моль}$ ($m=0$) & PASS \\
Барьер перескока $E_{\mathrm{hop}}^{(-1,0,+1)}$ & $s_0 E_{\mathrm{soft}}^{(m)} = \frac{1}{2}E_{\mathrm{soft}}^{(m)}$ & $3{,}0;\ 10{,}4;\ 20{,}8$ & $10{,}4\ \text{кДж/моль}$ ($m=0$) & PASS \\
Спайк мембран $\Delta V_{\mathrm{spike}}^{(-1,0,+1)}$ & $s_0 V_{\mathrm{soft}}^{(m)}$ ($V_{\mathrm{rest}} = -\frac{1}{3}V_{\mathrm{soft}}^{(m)}$) & $30{,}9;\ 108{,}0;\ 216{,}0$ & $107\ \text{мВ}$ ($V_{\mathrm{rest}}=-70$) & PASS \\
Фотосинтез $m=-1$ (Титан) & $[\lambda_{13}^{(-1)},\ \lambda_8^{(-1)}]$ ($\{13,8\}E_{\mathrm{soft}}^{(-1)}$) & $1{,}55\text{--}2{,}51\ \text{мкм}$ & ИК-окна $1{,}5$ и $2{,}0\text{--}2{,}5\ \text{мкм}$ & PASS \\
Фотосинтез $m=0$ (Земля) & $[\lambda_{13}^{(0)},\ \lambda_8^{(0)}]$ ($\{13,8\}E_{\mathrm{HB}}$) & $441{,}6\text{--}717{,}6\ \text{нм}$ & Хлорофилл (Соре--$P_{700}$) & PASS \\
Фотосинтез $m=+1$ (Венера) & $[\lambda_{13}^{(+1)},\ \lambda_8^{(+1)}]$ ($\{13,8\}E_{\mathrm{soft}}^{(+1)}$) & $220{,}8\text{--}358{,}8\ \text{нм}$ & УФ-поглотитель $220\text{--}360$ & PASS \\
Капсомеры $N_{\mathrm{shell}}(n)$ и $N_\varphi$ & $10n^2+2;\ 10\cdot 7+2;\ N_\varphi$ & $12,42,92,162;\ 72;\ 13$ & HBV ($42$), HSV ($162$) & PASS \\
Угол атаки Бюрги--Дюница $\theta_{\mathrm{BD}}$ & $\arccos(-1/d) = \arccos(-1/3)$ & $109{,}47^\circ$ & $107^\circ \pm 2^\circ$ ($K_3\to K_4$) & PASS \\
Катализ $E_a^{(0,2,3)}$ ($m=0$, кДж/моль) & $\{8,\ 3,\ 1/2\}E_{\mathrm{HB}}$ ($8=3+2+3$) & $166{,}7;\ 62{,}5;\ 10{,}4$ & $160;\ 60\text{--}65;\ 10\text{--}12$ & PASS \\
Ричардс и Трутон $\Sigma_{\mathrm{I\to II}},\ \Sigma_{\mathrm{II\to III}}$ & $9/7,\ 18/7;\ 21/2,\ 27/2$ & $1{,}29,\ 2{,}57;\ 10{,}5,\ 13{,}5$ & $1{,}2,\ 2{,}64;\ 10{,}5,\ 13{,}1$ & PASS \\
Упаковка $Z_c$ и Питцер $\omega(\mathrm{H_2O})$ & $7/24\ (\text{ГЦК});\ 3/13,\ 9/26$ & $0{,}292;\ 0{,}231,\ 0{,}346$ & $0{,}290;\ 0{,}229,\ 0{,}344$ & PASS \\
Индексы 3D-перехода $\nu,\ \gamma,\ \delta,\ \beta$ & $17/27;\ 68/55;\ 91/19$ & $0{,}630;\ 1{,}236;\ 4{,}79$ & $0{,}630;\ 1{,}237;\ 4{,}79$ & PASS \\
Распад $\mathrm{H_2O_2}$ ($-\Delta H_{\mathrm{rxn}}^{\circ}$, кДж/моль) & $5\,E_{\mathrm{HB}}\ (\text{г.});\ 4{,}5\,E_{\mathrm{HB}}\ (\text{ж.})$ & $104{,}2;\ 93{,}8$ & $105{,}5;\ 98{,}1$ (NIST) & PASS \\
Одеяло $\mathrm{N_2:O_2:Ar}$ и $T_{\mathrm{trop,surf}}$ & $63:17:1;\ \frac{7}{81},\frac{28}{243}T_{\mathrm{HB}}$ & $216{,}6;\ 288{,}8\ \text{К}$ & $216{,}7;\ 288{,}2\ \text{К}$ (ISA) & PASS \\
\bottomrule
\end{tabular}
\end{table}

\subsubsection*{Анализ, расчёт погрешности и выводы}
Неустранимая угловая щель пяти тетраэдров $\Delta\theta_{\mathrm{frust}} \approx 7{,}36^\circ$ объясняет принудительное скручивание биополимеров в апериодические спирали во всех трёх температурных окнах жизни ($m\in\{-1,0,+1\}$), а масштабный ряд $\{5,6,7\}\cdot 12\caKappaGeom a_0$ и двухбитный регистр~$\mathbb{Z}_4$ связывают в единую алгебру размеры мембран Титана, Земли и Венеры со спектром мутаций и законами Менделя.

\textbf{Вывод.}
Законы генетики и архитектура трёх планетарных типов жизни однозначно вытекают из геометрии кубооктаэдра~\latGCKshort.
\textbf{Статус записи:} верифицировано автоматическим гейтом \texttt{Bio\_frustration\_three\_life\_types}.
